% ===================================================================
% CHAPTER 5: METHODOLOGY
% ===================================================================
\chapter{Методологија} % Methodology
\label{chap:methodology}

Оваа глава ја опишува методологијата на експериментот следејќи ја логиката на процесниот тек (\emph{pipeline}): од податоците и нивната подготовка
(\hyperref[sec:data]{Дел~\ref*{sec:data}}), преку моделот и двата алгоритми за машинско заборавање што се споредуваат (\hyperref[sec:core_algorithm]{Дел~\ref*{sec:core_algorithm}}),
дефинирањето на контролираните експериментални услови и физичката мерна средина
(\hyperref[sec:exp_conditions]{Дел~\ref*{sec:exp_conditions}}), сè до рамката за
вреднување и статистичкиот метод
(\hyperref[sec:eval_framework]{Дел~\ref*{sec:eval_framework}}). Целта на
методологијата е споредба на потрошувачката на физички ресурси (време, енергија, топлина,
влезно-излезни операции на дискот) на двата пристапа за отстранување компромитирани податоци од
федеративен модел на работ на мрежата: наивно повторно тренирање и заборавање со SISA.

Целокупниот дизајн е претходно регистриран и замрзнат пред спроведувањето на
експерименталните извршувања, со што се спречува наменско избирање на резултатите по нивното добивање (\emph{p-hacking}).
Сите почетни вредности (\emph{seeds}) и поделбата на податоците се фиксни, така што секое
извршување е пресметковно детерминистичко и повторливо.

% ===================================================================
\section{Податоци} % Data
\label{sec:data}

\subsection{Извор и природа на податочното множество}
Експериментите се засноваат на податочното множество TON\_IoT
\cite{moustafa2021toniot}, објавено од UNSW, кое содржи мрежен сообраќај од експеримент со
уреди од Интернет на нештата (IoT), означен како бенигнен или напад. Се користат сите
23 CSV-датотеки со означен мрежен сообраќај (вкупно $\approx$3,3\,GB).
Задачата е формулирана како бинарна класификација - напад наспроти
бенигнен сообраќај - а не како повеќекласна класификација на нападите, што соодветствува
на улогата на систем за откривање упади (IDS).

\subsection{Предобработка}
Предобработката се изведува еднократно врз сите датотеки. Првично се задржуваат исклучиво нумеричките атрибути што се заеднички за сите датотеки (метаподатоците и колоните со ознаки се отстрануваат), со што се формира 15-димензионален вектор на обележја (features). Потоа, секоја партиција независно се стандардизира (со нулта средна вредност и единечна дисперзија), при што податоците во секоја партиција се делат на множества за тренирање и за валидација во сооднос 80/20, без вкрстено мешање меѓу јазлите. Бидејќи распределбата на податоците се врши на ниво на цели датотеки (а не на ниво на поединечни редови), добиените партиции по својата природа не се независни ниту идентично распределени (non-IID) според изворот, со што пореалистично се претставуваат хетерогените федеративни услови.

\subsection{Федеративна партиција и подмножество за експериментот}
За спроведување на експериментот, податоците се делат на четири партиции. Трите симулирани јазли (партициите 0--2) се распоредени како Docker-контејнери на истиот компјутер (лаптоп) на кој се извршува и централниот агрегатор, додека физичкиот уред Raspberry Pi ја содржи партицијата 3.

За да се изедначи придонесот на секој клиент во рамките на алгоритмот FedAvg, секој јазол се тренира врз подмножество од 1\,000\,000 реда. Одделно се издвојува чисто глобално множество за тестирање од 100\,000 реда, кое е стратификувано и дисјунктно од сите множества за тренирање (извлечено од валидациските региони врз кои никогаш не се тренира). Ова множество се користи за оценување на моделот од страната на домаќинот
(\hyperref[sec:eval_framework]{Дел~\ref*{sec:eval_framework}}).

Важна карактеристика на ова множество за тестирање претставува неговата изразена класна нерамнотежа: приближно 96,5\% напади наспроти 3,5\% бенигни примери. Оваа нерамнотежа има клучно влијание врз изборот на метриката за успешност и е детално образложена во \hyperref[sec:eval_framework]{Дел~\ref*{sec:eval_framework}}.

\subsection{Модел на закана и конструкција на труењето}
\label{subsec:poison}

Отстранувањето што е предмет на мерење е предизвикано од труење на податоците преку промена на етикетите (\emph{label-flipping}): дел од нападите ($y=1$) се преозначуваат како бенигни ($y=0$), со што се моделира извор што е компромитиран во моментот $\tau$. 
Труењето е детерминистичко во зависност од почетната вредност (\emph{seed}) и е просторно ограничено --- се однесува исклучиво на сегментите (\emph{slices}) 3 и 4 од шардот 1 (\emph{shard 1}) на компромитираниот јазол (партиција 3). Во овие сегменти се менуваат етикетите на сите напади (удел $\text{fraction}=1{,}0$), со што се претставува извор што е целосно компромитиран по моментот $\tau$. 
На овој начин се генерираат приближно 63\,000 компромитирани примери по јазол. Бидејќи распределбата на податоците во шардови и сегменти директно зависи од почетната вредност, оваа бројка благо варира при секое извршување. Заради транспарентност, точниот број на компромитирани примери секогаш е евидентиран во придружниот манифест. Кај две почетни вредности (50 и 51), подготовката на податоците е спроведена со удел од $0{,}5$, што резултира со околу 31\,500 компромитирани примери; ова отстапување и неговите последици се елаборирани во \ref{chap:limitations},.

Ова труење дефинира јасно специфицирано множество примери што мора да бидат заборавени. Во согласност со стандардите за машинско заборавање, отстранувањето на податоците не се спроведува поради нивното штетно влијание врз моделот, туку исклучиво затоа што се идентификувани како компромитирани (на пример, поради регулаторни барања). 

Дополнително, бидејќи алгоритмот FedAvg го разредува придонесот на компромитираниот клиент, труењето речиси и не влијае врз вкупната прецизност на глобалниот модел. Затоа, фокусот на ова истражување не е штетата од нападот, туку мерењето на трошокот за процесот на заборавање. 

Самото откривање на компромитираните податоци е надвор од опсегот на овој труд; наместо тоа, се претпоставува дека нивната локација е однапред позната и тие се експлицитно означени за отстранување.

% ===================================================================
\section{Модел и алгоритми за заборавање} % Core Algorithm/Model Description
\label{sec:core_algorithm}

\subsection{Модел за откривање упади}

Класификаторот претставува повеќеслоен перцептрон (MLP) составен од четири целосно поврзани слоеви со димензии $15\!\to\!128\!\to\!64\!\to\!32\!\to\!1$. Во секој скриен слој се применува комбинација од линеарна трансформација, активациска функција ReLU, пакетна нормализација (\emph{BatchNorm}) и испуштање (\emph{Dropout}, $p=0{,}3$), додека на излезниот неврон се применува сигмоидна активациска функција за пресметување на веројатноста дека даден мрежен тек претставува напад. Моделот се тренира преку функција за загуба базирана на бинарна вкрстена ентропија (\emph{BCELoss}), со користење на оптимизаторот Adam ($\text{lr}=0{,}001$).

\subsection{Федеративно учење}

Процесот на федеративно учење се извршува преку алгоритмот FedAvg, користејќи ја рамката Flower 1.29 во режим на распоредување (\emph{deployment}) со четири клиенти. Во секоја рунда, серверот ги испраќа хиперпараметрите, клиентите локално тренираат една епоха и ги враќаат тежините заедно со бројот на примери, по што серверот пресметува тежинска средна вредност. Применета е приспособена, отпорна варијанта на FedAvg: доколку во дадена рунда одговараат помалку клиенти од минимумот, се задржуваат тежините од претходната рунда, наместо да се агрегира нецелосно ажурирање. Ова спречува дивергенција на моделот при евентуални прекини на физичкиот јазол.

\subsection{Пристап A --- Наивно повторно тренирање}

Наивното заборавање служи како референтна основа (\emph{baseline}) за споредба со другите методи. Процесот е директен: компромитираните редови се отстрануваат од локалното множество, по што моделот се тренира целосно од почеток само врз преостанатите податоци. Ова ново тренирање трае 10 епохи (како и во првичната фаза) и користи нов, реиницијализиран оптимизатор. Иако овој пристап гарантира апсолутно заборавање (гарантирано со самата постапка), негов главен недостаток е неефикасноста, бидејќи бара целосно повторување на сите пресметки.

\subsection{Пристап Б --- Заборавање со SISA}

Методот SISA (\emph{Sharded, Isolated, Sliced, Aggregated}) \cite{bourtoule2021sisa} овозможува отстранување на податоците преку делумно повторно тренирање. Локалното множество се дели на $S=5$ шардови составени од по $R=5$ сегменти, користејќи пермутација базирана на почетната вредност (\emph{seed}). За секој шард се одржува посебен составен модел (\emph{constituent}) кој се тренира исклучиво врз тој шард. По обработката на секој сегмент, состојбата на моделот и на оптимизаторот се зачувува во контролна точка (\emph{checkpoint}). На крајот, глобалното ажурирање што се испраќа до алгоритмот FedAvg претставува параметарска средна вредност од сите составни модели.

При идентификување компромитиран сегмент $l^{*}$ во даден шард $s$, процесот на обновување (\emph{recovery}) се состои од:
(1) враќање (\emph{rollback}) на тој шард до контролната точка од последниот чист сегмент $(\text{рунда }1,\ \text{сегмент }l^{*}\!-\!1)$; и
(2) повторна обработка (\emph{replay}) исклучиво на засегнатите сегменти од тој шард. 

Незасегнатите шардови остануваат целосно недопрени. Во рамките на нашата поставка, труењето во сегментите 3--4 од шардот 1 резултира со враќање до $(\text{рунда }1,\ \text{сегмент }2)$ и повторно тренирање на 47 сегменти, за разлика од целосното повторно тренирање кое се бара кај наивното заборавање.

\paragraph{Отстапувања од класичниот SISA.} Овој труд воведува две важни промени во однос на оригиналниот метод SISA: (i) составните модели не ги преземаат глобалните тежини, што е неопходно за да може правилно да функционира враќањето до контролна точка; и (ii) кон серверот се испраќа параметарска средна вредност од составните модели, а не збирен резултат од нивните предвидувања. Оваа втора промена има директно влијание врз глобалната успешност на моделот, што е детално анализирано во \hyperref[sec:eval_framework]{Дел~\ref*{sec:eval_framework}} и во главата со резултати.

\paragraph{Опсег на заборавањето.} Двата пристапа гарантираат целосно локално заборавање --- обновениот модел на клиентот докажано не содржи никакви траги од труењето. Меѓутоа, тие нудат само приближно глобално заборавање, бидејќи клиентот не може да го избрише отровот што преку FedAvg веќе бил агрегиран во централниот сервер. Токму затоа, овие два метода не се споредуваат според тоа кој обезбедува поквалитетно заборавање, туку се споредуваат исклучиво според трошоците (ресурсите) потребни за да се постигне истиот краен резултат.

% ===================================================================
\section{Дефинирање на експерименталните услови}
\label{sec:exp_conditions}

\subsection{Главен фактор и протокол во четири фази}

Основниот независен фактор што се тестира во експериментот е методот на заборавање (Пристап A наспроти Пристап Б). Секое извршување на експериментот се одвива според идентичен протокол кој се состои од четири фази:

\begin{enumerate}
  \item \textbf{Фаза 1 --- Тренирање со труење:} 10 рунди федеративно тренирање со активно труење на компромитираниот јазол.
  \item \textbf{Фаза 2 --- Идентификација:} компромитираните примери експлицитно се означуваат за отстранување (самото откривање на нападот е надвор од опсегот на трудот).
  \item \textbf{Фаза 3 --- Обновување:} ова е главната фаза во која се вршат мерењата --- се применува наивно повторно тренирање или обновување преку SISA, извршено \emph{локално на уредот Raspberry Pi}.
  \item \textbf{Фаза 4 --- Повторно приклучување:} 5 дополнителни рунди кои продолжуваат од глобалниот модел (во кој претходно веќе е агрегирано влијанието од труењето). Бидејќи компромитираните податоци на клиентот сега се веќе отстранети, во оваа фаза се мери влијанието врз глобалната успешност на моделот.
\end{enumerate}

\subsection{Физичка експериментална околина}
\label{subsec:testbed}

Слика~\ref{fig:testbed} ја прикажува експериментална околина во која се извршени сите мерења. Ограничувањата на хардверот се намерни со цел реално да се симулира работната околина на еден типичен IoT-јазол. Експериментот се спроведува со следните хардверски компоненти и системски услови:

\begin{itemize}
  \item \textbf{Пресметувачка моќ и температура:} Raspberry Pi 5 (4\,GB) со намерно исклучен вентилатор за ладење, што предизвикува термичко придушување (\emph{thermal throttling}) при 85\,\textdegree C.
  \item \textbf{Влезно-излезни операции на дискот (I/O):} 128\,GB SD-картичка од класа A1 (ограничена на $\approx$500 случајни запишувања во секунда). Ова претставува критично тесно грло за методот SISA поради потребата од често зачувување на контролните точки.
  \item \textbf{Напојување:} надворешен мерач FNIRSI FNB58, континуирано поврзан на напојувањето на јазолот за време на сите мерења. Се користи кабел со поддршка за струја од 5\,A за да се спречи пад на напонот, што би предизвикало несакано забавување на системот.
  \item \textbf{Изолација на системскиот шум:} виртуелната меморија (\emph{swap}) е целосно оневозможена, а телеметријата се запишува исклучиво во RAM-меморијата. На овој начин, метриката за I/O операциите на дискот го одразува само вистинското работно оптоварување од алгоритмите.
\end{itemize}

% -------------------------------------------------------------------
% TODO (СЛИКА НА ХАРДВЕР): Овде вметни ФОТОГРАФИЈА од физичката поставка ---
% Raspberry Pi 5 со инлајн поврзан FNB58 мерач и кабелажата за напојување/Ethernet.
% (Логичката/системската топологија на Flower контејнерите оди во поглавјето
%  „Архитектура“, НЕ тука, за да не се дуплира.)
% -------------------------------------------------------------------
\begin{figure}[htbp]
  \centering
  % \includegraphics[width=0.75\textwidth]{figures/hardware_testbed.jpg}
  \fbox{\parbox[c][0.28\textheight][c]{0.75\textwidth}{\centering
    \textit{[Место за фотографија од физичката експериментална околина:\\
    Raspberry Pi 5 (без вентилатор) + мерач FNB58 во серија + кабли]}}}
  \caption{Физичка експериментална околина: Raspberry Pi 5 со намерно исклучен вентилатор,
  FNIRSI FNB58 мерач поставен сериски (inline) на напојувањето, како и кабли за
  напојување и мрежа. Мерачот е присутен при сите извршувања, така што неговиот
  режиски трошок претставува константа, а не збунувачки фактор (\emph{confound}).}
  \label{fig:testbed}
\end{figure}

\subsection{Фиксни експериментални параметри}

Табела~\ref{tab:fixed-params} ги сумира параметрите кои се дефинирани по пилот-тестирањето и кои остануваат фиксни за време на сите главни експериментални извршувања.

\begin{table}[htpb]
  \centering
  \caption{Преглед на фиксните експериментални параметри.}
  \label{tab:fixed-params}
  \begin{tabular}{ll}
    \toprule
    \textbf{Параметар} & \textbf{Вредност} \\
    \midrule
    Експериментални пристапи & A: Наивно повторно тренирање · B: SISA \\
    SISA конфигурација & $S=5$ шардови $\times$ $R=5$ сегменти \\
    Труење на податоци & Напад претворен во бениген, удел $1{,}0$, сегменти 3--4 во шард 1 \\
    Мрежна топологија & 4 клиенти (3 симулирани + 1 физички Pi преку LAN) \\
    Податоци по клиент & 1\,000\,000 примери (стратификувано подмножество) \\
    Множество за тестирање & 100\,000 примери, дисјунктно, $\approx$96,5\% напади \\
    Број на рунди & Фаза 1: 10 рунди · Фаза 4: 5 рунди \\
    Времетраење (Наивен метод) & 10 епохи врз преостанатите чисти податоци \\
    Почетни вредности (\emph{seeds}) & $N=10$ спарени опсервации (42--51) \\
    Редослед на извршување & Балансиран (A/B) според почетните вредности \\
    Статистички тестови & Wilcoxon, Cliff's $\delta$, 95\% интервал на доверба (\emph{bootstrap}) \\
    \bottomrule
  \end{tabular}
  \par\medskip
  \footnotesize\raggedright
  \emph{Забелешка:} За почетните вредности 50 и 51, подготовката на податоците е спроведена со удел на труење од $0{,}5$; влијанието на ова отстапување е разгледано во \hyperref[chap:limitations]{Дел~\ref*{chap:limitations}}.
\end{table}

\subsection{Контрола на надворешни фактори}

За да се обезбеди изолирана и објективна споредба помеѓу двата метода, применети се следните механизми:

\begin{itemize}
  \item \textbf{Термичка стабилизација:} Пред секое мерење, системот паузира сѐ додека температурата на процесорот не се стабилизира (промена од најмногу $2\,\textdegree$C во период од 120\,s). Бидејќи Raspberry Pi без вентилатор не може целосно да се излади, извршувањата започнуваат од стабилна термичка состојба, а не од конкретна температурна вредност.
  \item \textbf{Балансиран редослед:} За редоследот на извршување да не влијае врз резултатите, тој рамномерно се ротира (половина од експериментите започнуваат со наивниот пристап, а половина со SISA).
  \item \textbf{Софтверски детерминизам:} Почетните вредности (\emph{seeds}) за библиотеките (\emph{PyTorch, NumPy}) и за поделбата на податоците се строго фиксирани. Ова гарантира дека пресметките се целосно повторливи (иста вредност на загуба при секое извршување), додека времето на извршување природно варира како физичко мерење.
  \item \textbf{Поврзан дизајн (\emph{paired design}):} Двата метода се извршуваат на истиот физички уред под исти услови. Така, за секоја почетна вредност се добива пар поврзани мерења, што овозможува прецизна статистичка споредба.
\end{itemize}

Дополнително, со цел да се симулира хетерогена околина, на трите симулирани клиенти во сите извршувања им се доделени различни процесорски ограничувања (3,0 / 2,0 / 1,5 јадра). Како дополнителна проверка на робусноста, еден пар извршувања (двата метода за иста почетна вредност) се спроведува во услови на реалистична WAN мрежна латентност (\texttt{toxiproxy}, 40\,ms $\pm$ 20\,ms на Fleet API). Целта е да се потврди валидноста на резултатите и надвор од идеалните услови на една локална мрежа (LAN).

% ===================================================================
\section{Методологија на евалуација}
\label{sec:eval_framework}

Евалуацијата се заснова на шест претходно дефинирани хипотези. Во продолжение се опишани зависните променливи и соодветните мерни инструменти.

\subsection{Показатели за потрошувачка на физичките ресурси (H1--H4, H6)}
\begin{itemize}
  \item \textbf{H1 --- Време до обновување (TTR):} Времетраењето на процесот на обновување (Фаза 3) во секунди, претставено како вкупно изминато време (\emph{wall-clock time}) отчитано директно од модулот.
  \item \textbf{H2 --- Потрошувачка на енергија при обновувањето:} Вкупната енергија потрошена во текот на Фаза 3, пресметана како трапезоиден интеграл $\int W\,dt$. Податоците се прибираат преку мерачот FNB58 ($\approx$100 мерења во секунда) со усогласени временски ознаки. Прикажана е исклучиво нето-вредноста (дополнителната потрошувачка), која се добива кога од вкупната ќе се одземе потрошувачката на уредот во состојба на мирување.
  \item \textbf{H3 --- Топлотно оптоварување:} Основен показател е времетраењето на активното температурно ограничување (во секунди) за време на Фаза 3. Вредноста се отчитува од системските регистри на процесорот исклучиво преку тековно активните битови, за да се спречат неточни резултати од претходните состојби. Најниската работна фреквенција се бележи само информативно бидејќи не покажува статистички значајна разлика меѓу пристапите (паѓа само кај дел од извршувањата), додека највисоката температура не се зема предвид бидејќи брзо се стабилизира на горната граница од $\approx$87\,\textdegree C.
  \item \textbf{H4 --- Влезно-излезни (I/O) операции на дискот:} Вкупниот обем на податоци запишани на SD-картичката (во килобајти) и просечното време на чекање на дискот (\emph{iowait}) за време на Фаза 3. Овој показател има истражувачки карактер и се пријавува објективно, дури и кога неговото влијание е занемарливо.
  \item \textbf{H6 --- Режиски трошоци на SISA:} Дополнителното време по рунда и обемот на податоци запишани во контролните точки во текот на првичното тренирање (Фаза 1) — особина специфична исклучиво за SISA. Вредностите се пресметуваат преку уникатните запишувања по сегмент, при што податоците од останатите фази не се земаат предвид. Вкупниот ресурсен трошок претставува збир од режискиот трошок во Фаза 1 и енергијата потрошена за самото обновување.
\end{itemize}

Скриптата за телеметрија на уредот Raspberry Pi прибира податоци со зачестеност од 1\,Hz (односно, создава по еден запис секоја секунда) и ги бележи следниве вредности: температурата, работната фреквенција, температурното ограничување, времето на чекање на дискот (\emph{iowait}) и обемот на податоци запишани на SD-картичката. Секој запис соодветно се обележува со тековната фаза од експериментот. Истовремено, главниот компјутер (домаќинот) ги презема податоците од мерачот FNB58 директно преку USB-врска.

\subsection{Показатели за успешност на моделот (H5)}
\label{subsec:utility-metrics}

Хипотезата H5 проверува дали обновениот модел ја задржува својата успешност. Бидејќи глобалното множество за тестирање има екстремна класна нерамнотежа (96,5\% напади), стандардните показатели (одѕив и мерка F1) даваат тривијални и залажувачки резултати. Затоа, успешноста се оценува преку показатели отпорни на класна нерамнотежа, изведени од матрицата на грешки:

\begin{itemize}
  \item \textbf{Специфичност:} уделот на точно препознаени бенигни примери;
  \item \textbf{Избалансирана точност:} $=(\text{одѕив} + \text{специфичност}) / 2$ (вредност од $0{,}50$ означува случаен избор);
  \item \textbf{Метјуов коефициент на корелација (MCC):} (вредност од $0$ означува случаен избор);
  \item \textbf{ROC-AUC:} способност за рангирање независна од прагот на одлучување.
\end{itemize}


Одѕивот (\emph{recall}), прецизноста (\emph{precision}) и мерката F1 се пријавуваат само заради целовитост на прегледот, додека заклучоците се темелат исклучиво на избалансираната точност, специфичноста и MCC.

\paragraph{Отстапување од првичниот протокол.} Првичниот план за H5 се потпираше исклучиво на толеранција од $\Delta F1 \leq 2$ процентни поени во однос на основниот модел. Сепак, по завршувањето на експериментите беше воочено дека мерката F1 е несоодветна за вака нерамнотежено множество. Затоа, \emph{post hoc} беше воведена евалуацијата со избалансираните показатели. Ова е свесно отстапување од почетниот дизајн, а оригиналните F1-вредности остануваат достапни во резултатите (\texttt{results\_phase\{1,4\}.json}) заради транспарентност.

Евалуацијата (преку скриптата \texttt{analysis/reeval\_utility.py}) се врши врз зачуваните модели од Фаза 1 и Фаза 4. Покрај мерењето при стандардниот праг на одлучување ($0{,}5$), за секој модел се пресметува и максимално достижната избалансирана точност низ сите можни прагови. Разликата помеѓу овие две вредности ја изолира загубата во успешноста која се должи исклучиво на несоодветна калибрација. Резултатите се анализирани во \hyperref[chap:results]{Глава~\ref*{chap:results}}.

\subsection{Статистички метод}

Експерименталниот дизајн се заснова на пар поврзани мерења за секоја почетна вредност (\emph{seed}). Поради тоа, за анализата се користат непараметриски статистички методи: Вилкоксоновиот тест на рангирани знаци (\emph{Wilcoxon signed-rank test}), коефициентот Cliff's $\delta$ за мерење на големината на ефектот и 95\% интервал на доверба (\emph{bootstrap}) за медијаната на разликата меѓу паровите.

Бидејќи примерокот е мал ($N=10$), минималната можна $p$-вредност при двостран тест изнесува $0{,}0020$. Поради ова ограничување, заклучоците се темелат првенствено на големината на ефектот и на интервалите на доверба, наместо исклучиво на статистичката значајност ($p$-вредноста). Сите сурови податоци се јавно достапни.

Сумарните статистики за потрошувачката на физичките ресурси се дадени во Табела~\ref{tab:paired-cost}, додека Табела~\ref{tab:utility} ги прикажува резултатите за успешноста преку избалансираните показатели. Детална анализа на овие податоци следи во главата со резултати.